\chapter{The Alternation Principle}
\label{ch:ii23}

The alternation principle is the structural certificate for best uniform approximation. It converts an optimization problem into a finite pattern of equal-magnitude alternating errors.

\section*{Learning goals}
After this chapter, the reader should be able to:
\begin{itemize}
\item state the alternation criterion with the correct number and ordering of alternating extrema;
\item use an alternation pattern to certify that a polynomial approximant is minimax;
\item construct a perturbation argument showing that insufficient alternation cannot be optimal;
\item identify the error function and its extreme points in explicit approximation problems;
\item distinguish necessary from sufficient parts of the alternation principle under the chapter's hypotheses;
\item apply the principle to low-degree examples without solving a larger optimization problem directly;
\end{itemize}
\section*{Conceptual roadmap}
\[
\boxed{\text{definition}\;\longrightarrow\;\text{estimate}\;\longrightarrow\;\text{theorem}\;\longrightarrow\;\text{application}.}
\]

\section{Error function}
For approximant $p$, study $e=f-p$.

\section{Extreme points}
Minimax optimality is witnessed where $|e|$ reaches its sup norm.

\section{Alternation pattern}
Signs alternate at successive extremal points.

\section{Necessity}
Without enough alternation one can construct a perturbation reducing the maximum error.

\section{Sufficiency}
Enough alternating extrema block every lower-degree perturbation by root counting.

\section{Uniqueness}
Two best approximants would force their difference to have too many zeros.

\section{Remez idea}
Algorithms update candidate alternation points until equal error levels stabilize.

\section{Approximation certificates}
Alternation gives a verifiable finite certificate of global optimality.

\section{Core structural results}

\begin{theorem}[Chebyshev alternation theorem]\label{thm:ii23-01}
A polynomial $p$ of degree at most $n$ is best uniform approximation to continuous $f$ on $[a,b]$ iff the error attains its maximum magnitude at at least $n+2$ ordered points with alternating signs.
\begin{proof}
Sufficiency follows because any better competitor would force the difference polynomial to change sign between each adjacent pair, producing more roots than its degree. Necessity follows from a separation/perturbation argument when no such alternating set exists.
\end{proof}
\end{theorem}

\begin{theorem}[Uniqueness]\label{thm:ii23-02}
The best degree-$n$ uniform polynomial approximant is unique.
\begin{proof}
If two best approximants existed, alternation for one and the competitor inequalities would force their difference to have at least $n+1$ zeros despite degree at most $n$.
\end{proof}
\end{theorem}

\begin{theorem}[Finite certificate]\label{thm:ii23-03}
Once $n+2$ alternating equal-magnitude error points are verified, global minimax optimality follows.
\begin{proof}
This is the sufficiency direction of the alternation theorem.
\end{proof}
\end{theorem}

\section{Worked examples}

\begin{example}[Constant approximation]\label{ex:ii23-worked-01}
How many alternating points certify a best constant approximant? Two points with errors $+E$ and $-E$; the best constant centers the range between maximum and minimum values.
\end{example}

\begin{example}[Best constant formula]\label{ex:ii23-worked-02}
For continuous $f$, find the best constant in sup norm. It is $(M+m)/2$, where $M=\max f$ and $m=\min f$, with error $(M-m)/2$.
\end{example}

\begin{example}[Degree-one certificate]\label{ex:ii23-worked-03}
How many alternating extrema certify a best line? Three.
\end{example}

\begin{example}[Why equal magnitude?]\label{ex:ii23-worked-04}
Why must active extrema in the alternation pattern have the same magnitude? If one were strictly smaller, the true sup norm would be determined elsewhere and the finite optimality certificate would not be balanced.
\end{example}

\section{Solved dossiers}
The dossiers are longer solved problems. Legacy mappings guide topic selection where explicit retained source blocks exist; otherwise fresh canonical problems complete the chapter.

\begin{problem}[Constant approximation]\label{prob:ii23-dossier-01}
How many alternating points certify a best constant approximant?
\end{problem}
\begin{solution}
Two points with errors $+E$ and $-E$; the best constant centers the range between maximum and minimum values.
\end{solution}

\begin{problem}[Best constant formula]\label{prob:ii23-dossier-02}
For continuous $f$, find the best constant in sup norm.
\end{problem}
\begin{solution}
It is $(M+m)/2$, where $M=\max f$ and $m=\min f$, with error $(M-m)/2$.
\end{solution}

\begin{problem}[Degree-one certificate]\label{prob:ii23-dossier-03}
How many alternating extrema certify a best line?
\end{problem}
\begin{solution}
Three.
\end{solution}

\begin{problem}[Why equal magnitude?]\label{prob:ii23-dossier-04}
Why must active extrema in the alternation pattern have the same magnitude?
\end{problem}
\begin{solution}
If one were strictly smaller, the true sup norm would be determined elsewhere and the finite optimality certificate would not be balanced.
\end{solution}

\begin{problem}[Root-count sufficiency]\label{prob:ii23-dossier-05}
Explain the sign-change argument.
\end{problem}
\begin{solution}
A uniformly better competitor must lie closer to $f$ at each alternating extremum, forcing the difference of the two approximants to alternate sign between consecutive extremal points and thus have too many roots.
\end{solution}

\begin{problem}[Uniqueness]\label{prob:ii23-dossier-06}
Why can two different minimax polynomials not share the same alternation certificate?
\end{problem}
\begin{solution}
Their difference would need to vanish or change sign in too many intervals for its degree.
\end{solution}

\begin{problem}[Remez iteration]\label{prob:ii23-dossier-07}
What does a Remez step update?
\end{problem}
\begin{solution}
It solves for an approximant and common error level on a chosen alternating set, then replaces points using extrema of the new error.
\end{solution}

\begin{problem}[Error centering]\label{prob:ii23-dossier-08}
For constant approximation, why center the maximum and minimum?
\end{problem}
\begin{solution}
Any shift away from the midpoint increases the larger of the two endpoint range deviations.
\end{solution}

\begin{problem}[Monotone target]\label{prob:ii23-dossier-09}
If $f$ is monotone, where are extreme errors for the best constant?
\end{problem}
\begin{solution}
At points attaining the minimum and maximum of $f$.
\end{solution}

\begin{problem}[Certificate versus search]\label{prob:ii23-dossier-10}
Why is alternation valuable computationally?
\end{problem}
\begin{solution}
Searching for the optimum is global, but once a candidate displays the required alternation, optimality can be certified finitely.
\end{solution}

\begin{problem}[Generalized spaces]\label{prob:ii23-dossier-11}
What property of polynomials enables alternation arguments?
\end{problem}
\begin{solution}
A nonzero degree-$n$ polynomial cannot have more than $n$ roots; more generally Chebyshev systems have controlled zero counts.
\end{solution}

\begin{problem}[Alternation synthesis]\label{prob:ii23-dossier-12}
What does the theorem reveal conceptually?
\end{problem}
\begin{solution}
Best worst-case approximations distribute their unavoidable error evenly with alternating sign rather than allowing one-sided or localized excess.
\end{solution}

\section{Exercises with complete solutions}

\begin{exercise}\label{exr:ii23-01}
Best constant to $f$ uses which values?
\end{exercise}
\begin{hint}
Maximum and minimum.
\end{hint}
\begin{solution}
Their midpoint.
\end{solution}

\begin{exercise}\label{exr:ii23-02}
How many points for degree $n$?
\end{exercise}
\begin{hint}
$n+2$.
\end{hint}
\begin{solution}
At least $n+2$.
\end{solution}

\begin{exercise}\label{exr:ii23-03}
What alternates?
\end{exercise}
\begin{hint}
Error sign.
\end{hint}
\begin{solution}
Signs of extremal error.
\end{solution}

\begin{exercise}\label{exr:ii23-04}
What is equal?
\end{exercise}
\begin{hint}
Error magnitude.
\end{hint}
\begin{solution}
The sup error magnitude.
\end{solution}

\begin{exercise}\label{exr:ii23-05}
What proves sufficiency?
\end{exercise}
\begin{hint}
Root count.
\end{hint}
\begin{solution}
Sign changes and root counting.
\end{solution}

\begin{exercise}\label{exr:ii23-06}
What algorithm uses alternation?
\end{exercise}
\begin{hint}
Remez.
\end{hint}
\begin{solution}
Remez algorithm.
\end{solution}

\begin{exercise}\label{exr:ii23-07}
Is minimax best approximant unique?
\end{exercise}
\begin{hint}
For polynomials, yes.
\end{hint}
\begin{solution}
Yes.
\end{solution}

\begin{exercise}\label{exr:ii23-08}
What norm is involved?
\end{exercise}
\begin{hint}
Sup norm.
\end{hint}
\begin{solution}
Uniform norm.
\end{solution}

\section*{Chapter summary}
The chapter's tools now form part of the canonical Volume II analysis toolkit and will be used in later approximation, fixed-point, and convergence arguments.
